\author{Jessica Schmidt}
\title{Intensitätspartikeln im Deutschen}
\subtitle{Zur Semantik von deskriptiver und expressiver Verstärkung}
\BackBody{Dieses Buch befasst sich mit dem in der theoretischen linguistischen Literatur stark diskutierten Phänomen der lexikalischen Adjektiv-Intensivierung. Dabei wird ein graduierbares Adjektiv – z. B. \textit{schön} – durch eine Intensitätspartikel modifiziert. Durch die Modifikation erfährt das Adjektiv eine Bedeutungsverstärkung, die sich in einer gesteigerten Intensität und einer höheren Position auf einer zugehörigen mentalen Skala niederschlägt.

\ea
    \ea     Der Strand ist \textit{∅ schön}.    \hfill         Positiv
    \ex     Der Strand ist \textit{\textbf{sehr} schön}. \hfill       Deskriptiv
    \ex     Der Strand ist \textit{\textbf{mega} schön}. \hfill     Expressiv
    \z
\z


Im Fokus steht die experimentelle Untersuchung von deskriptiven und expressiven Intensitätspartikeln. „Expressive Intensivierer“ (Gutzmann \& Turgay 2012) wie \textit{mega} und \textit{super} treten vor allem in der Umgangssprache auf, meist in Gestalt von Adjektiven, und erfüllen eine ähnliche Funktion wie „deskriptive Intensivierer“ (ebd.), etwa die deutsche Standard-Intensitätspartikel \textit{sehr}. In beiden Fällen wird im Vergleich zum Positiv ein höheres Ausmaß der bezeichneten Eigenschaft ausgedrückt; bei expressiven Intensivierern wird dieser Grad jedoch – so die in der Literatur vorherrschende Annahme – intuitiv als stärker empfunden. Zudem offenbaren expressive Partikeln aufgrund einer emotional-expressiven Zusatzkomponente mutmaßlich auch die persönliche Einstellung der sprechenden Person zum Sachverhalt, die nicht Teil der Proposition ist.

Obwohl die Erforschung von Expressivität in der formal orientierten Sprachwissenschaft in den letzten drei Jahrzehnten erheblich an Dynamik gewonnen hat, ist die Gültigkeit dieser grundlegenden Annahmen bislang nicht empirisch abgesichert. Die noch bestehende Lücke an experimentellen Untersuchungen zu schließen, ist eines der zentralen Ziele dieser Arbeit. Aufbauend auf der aktuellen Forschungslage unternimmt sie den ersten Versuch, durch eine breit angelegte Untersuchung neue Erkenntnisse über den semantischen Unterschied zwischen deskriptiven und expressiven Intensitätspartikeln zu gewinnen. Hierfür wird auf eine systematisch aufeinander aufbauende Kombination aus verschiedenen linguistischen Methoden wie Korpusanalyse (Empirieteil I) und Rating-Experimente (Empirieteil II) zurückgegriffen, mithilfe derer die aus der Literatur hergeleiteten Forschungshypothesen erstmals überprüft werden.
}
\renewcommand{\lsSeries}{ogl}
\renewcommand{\lsSeriesNumber}{17}
\renewcommand{\lsID}{583}
\lsCoverTitleSizes{48pt}{16.5mm}
%\renewcommand{\lsImpressumExtra}{Dieses Buch ist eine überarbeitete Fassung der Dissertation der Autorin, die sie im
%Dezember 2022 unter dem Titel \textit{Ist} sau cool \textit{cooler als} sehr cool\textit{? Empirische
%%Untersuchungen zur Semantik von deskriptiven und expressiven Ausdrücken am
%Beispiel deutscher Intensitätspartikeln} an der Universität des Saarlandes in Saarbrücken
%eingereicht und im Juli 2023 erfolgreich verteidigt hat.}

\renewcommand{\lsISBNdigital}{978-3-96110-563-2}
\renewcommand{\lsISBNhardcover}{978-3-98554-183-6}
\BookDOI{10.5281/zenodo.18999426}
\proofreader{Adriana Rosalina Galván Torres,
Jean Nitzke,
Katja Politt,
Laura Grestenberger,
Ludger Paschen,
Nicole Benker,
Tom Bossuyt,
Vicky Reiter,
Yvonne Treis
}
\typesetter{Jessica Schmidt, Bruno Behling, Hannah Schleupner}
